\documentclass[11pt,a4paper]{article}
\usepackage[a4paper,top=3.6cm,bottom=2.4cm,left=1.8cm,right=1.8cm,headheight=44pt,headsep=14pt]{geometry}
\usepackage{amsmath,amssymb,mathtools}
\usepackage{graphicx}
\usepackage{xcolor}
\usepackage{array,tabularx,longtable,booktabs}
\usepackage{fancyhdr}
\usepackage{lastpage}
\usepackage{needspace}
\usepackage{multicol}
\setlength{\columnsep}{0.9cm}\setlength{\columnseprule}{0.4pt}
\setlength{\parindent}{0pt}
\setlength{\emergencystretch}{3em}
\definecolor{schoolblue}{RGB}{20,50,120}
\pagestyle{fancy}
\fancyhf{}
\renewcommand{\headrulewidth}{0.8pt}
\renewcommand{\footrulewidth}{0.5pt}
\fancyhead[C]{\textcolor{schoolblue}{\Large\bfseries SAGAR PUBLIC SCHOOL}\\[2pt]
  \small\textbf{Integrated Programme}\ \ $\mid$\ \ \textbf{Topic: Electrostatics -- Coulomb's Law, Electric Field \& Field Lines}}
\fancyfoot[L]{\small Sagar Public School -- Integrated Programme}
\fancyfoot[C]{\small Page \thepage\ of \pageref{LastPage}}
\fancyfoot[R]{\small Practice Sheet 1\ $\bullet$\ Electrostatics}
\fancypagestyle{plain}{\fancyhf{}\renewcommand{\headrulewidth}{0.8pt}\renewcommand{\footrulewidth}{0.5pt}%
  \fancyhead[C]{\textcolor{schoolblue}{\Large\bfseries SAGAR PUBLIC SCHOOL}\\[2pt]\small\textbf{Integrated Programme}\ \ $\mid$\ \ \textbf{Topic: Electrostatics -- Coulomb's Law, Electric Field \& Field Lines}}%
  \fancyfoot[L]{\small Sagar Public School -- Integrated Programme}\fancyfoot[C]{\small Page \thepage\ of \pageref{LastPage}}\fancyfoot[R]{\small Practice Sheet 1\ $\bullet$\ Electrostatics}}
\newcounter{qn}
\newcommand{\Q}{\needspace{5\baselineskip}\refstepcounter{qn}\par\medskip\noindent\textbf{Q.\theqn}\hspace{0.5em}}
\newcommand{\Qf}{\needspace{7.5cm}\refstepcounter{qn}\par\medskip\noindent\textbf{Q.\theqn}\hspace{0.5em}}
\newcommand{\lab}[1]{(\ifcase#1\or A\or B\or C\or D\fi)}
\newcommand{\optsfour}[4]{\par\smallskip\noindent\begin{tabularx}{\linewidth}{@{}*{4}{>{\raggedright\arraybackslash}X}@{}}
  \lab1\ #1 & \lab2\ #2 & \lab3\ #3 & \lab4\ #4\end{tabularx}\par}
\newcommand{\optstwo}[4]{\par\smallskip\noindent\begin{tabularx}{\linewidth}{@{}*{2}{>{\raggedright\arraybackslash}X}@{}}
  \lab1\ #1 & \lab2\ #2 \\[4pt] \lab3\ #3 & \lab4\ #4\end{tabularx}\par}
\newcommand{\optsv}[4]{\par\smallskip\begin{list}{}{\setlength{\leftmargin}{1.7em}\setlength{\itemsep}{1pt}\setlength{\topsep}{0pt}\setlength{\parsep}{0pt}}
  \item[\lab1] #1 \item[\lab2] #2 \item[\lab3] #3 \item[\lab4] #4\end{list}}
\newcommand{\fig}[1]{\par\smallskip\begin{center}\includegraphics[width=\linewidth,height=4.6cm,keepaspectratio]{figs/#1.png}\end{center}\par}
\newcommand{\etag}[1]{\ {\footnotesize\textit{[#1]}}}
\begin{document}
\begin{center}
{\Large\bfseries PRACTICE SHEET 1}\\[3pt]
{\normalsize\bfseries Coulomb's Law \& Electrostatic Force $\bullet$ Coulomb's Law \& Charge Configurations $\bullet$ Electric Field due to Charge Distributions $\bullet$ Electric Field \& Field Lines}
\end{center}
\vspace{2pt}
\noindent\fbox{\parbox{\dimexpr\linewidth-2\fboxsep-2\fboxrule\relax}{%
\textbf{Time Allowed: 60 Mins}\hfill \textbf{Total Questions: 20}\\[5pt]
Name: \underline{\hspace{6.3cm}}\hfill Class/Sec: \underline{\hspace{2.2cm}}\hfill Date: \underline{\hspace{2.6cm}}\\[5pt]
{\small\textbf{Instructions:} Attempt all questions. Each question carries one correct option unless marked \emph{one or more than one option is correct}. Questions without options are numerical-value type: write the answer in the blank. The tag after each question shows its source examination. A self-assessment page follows; the answer key is provided as a separate sheet.}}}
\bigskip
\begin{multicols}{2}
\Q In hydrogen like system the ratio of coulombian force and gravitational force between an electron and a proton is in the order of:\etag{JEE Main 2024, 5 Apr, Mor.}
\optsfour{$10^{19}$}{$10^{39}$}{$10^{29}$}{$10^{36}$}

\Q Force between two point charges $q_1$ and $q_2$ placed in vacuum at $r$ cm apart is $F$. Force between them when placed in a medium having dielectric constant $K=5$ at $r/5$ cm apart will be:\etag{JEE Main 2024, 31 Jan, Eve.}
\optsfour{$5F$}{$25F$}{$F/5$}{$F/25$}

\Q If two charges $q_1$ and $q_2$ are separated with distance $d$ and placed in a medium of dielectric constant $K$. What will be the equivalent distance between charges in air for the same electrostatic force?\etag{JEE Main 2023, 24 Jan, Mor.}
\optstwo{$d\sqrt{K}$}{$1.5\,d\sqrt{K}$}{$K\sqrt{d}$}{$2d\sqrt{K}$}

\Q A charge of $4\,\mu$C is to be divided into two. The distance between the two divided charges is constant. The magnitude of the divided charges so that the force between them is maximum, will be:\etag{JEE Main 2022, 27 Jul, Eve.}
\optstwo{$1\,\mu$C and $3\,\mu$C}{$2\,\mu$C and $2\,\mu$C}{$0$ and $4\,\mu$C}{$1.5\,\mu$C and $2.5\,\mu$C}

\Q Three charges $+Q$, $q$, $+Q$ are placed respectively, at distance $0$, $d/2$ and $d$ from the origin, on the $x$-axis. If the net force experienced by $+Q$, placed at $x=0$, is zero, then value of $q$ is:\etag{JEE Main 2019, 9 Jan, Mor.}
\optsfour{$-\dfrac{Q}{4}$}{$+\dfrac{Q}{2}$}{$+\dfrac{Q}{4}$}{$-\dfrac{Q}{2}$}

\Qf The three charges $q/2$, $q$ and $q/2$ are placed at the corners A, B and C of a square of side $a$ as shown in figure. The magnitude of electric field ($E$) at the corner D of the square, is:\etag{JEE Main 2022, 28 Jun, Mor.}
\fig{f105}
\optsv{$\dfrac{q}{4\pi\epsilon_0a^2}\left(\dfrac{1}{\sqrt2}+\dfrac12\right)$}{$\dfrac{q}{4\pi\epsilon_0a^2}\left(1+\dfrac{1}{\sqrt2}\right)$}{$\dfrac{q}{4\pi\epsilon_0a^2}\left(1-\dfrac{1}{\sqrt2}\right)$}{$\dfrac{q}{4\pi\epsilon_0a^2}\left(\dfrac{1}{\sqrt2}-\dfrac12\right)$}

\Q Ten charges are placed on the circumference of a circle of radius $R$ with constant angular separation between successive charges. Alternate charges 1, 3, 5, 7, 9 have charge $(+q)$ each, while 2, 4, 6, 8, 10 have charge $(-q)$ each. The potential $V$ and the electric field $E$ at the centre of the circle are respectively. (Take $V=0$ at infinity)\etag{JEE Main 2020, 5 Sep, Eve.}
\optsv{$V=0;\ E=0$}{$V=\dfrac{10q}{4\pi\epsilon_0R};\ E=\dfrac{10q}{4\pi\epsilon_0R^2}$}{$V=\dfrac{10q}{4\pi\epsilon_0R};\ E=0$}{$V=0;\ E=\dfrac{10q}{4\pi\epsilon_0R^2}$}

\Q Four point charges $-q$, $+q$, $+q$ and $-q$ are placed on $y$-axis at $y=-2d$, $y=-d$, $y=+d$ and $y=+2d$, respectively. The magnitude of the electric field $E$ at a point on the $x$-axis at $x=D$, with $D\gg d$, will behave as:\etag{JEE Main 2019, 9 Apr, Eve.}
\optstwo{$E\propto\dfrac{1}{D^3}$}{$E\propto\dfrac{1}{D}$}{$E\propto\dfrac{1}{D^4}$}{$E\propto\dfrac{1}{D^2}$}

\Q Four charges equal to $-Q$ are placed at the four corners of a square and a charge $q$ is at its center. If the system is in equilibrium the value of $q$ is\etag{AIEEE 2004}
\optstwo{$-\dfrac{Q}{2}\left(1+2\sqrt2\right)$}{$\dfrac{Q}{4}\left(1+2\sqrt2\right)$}{$-\dfrac{Q}{4}\left(1+2\sqrt2\right)$}{$\dfrac{Q}{2}\left(1+2\sqrt2\right)$}

\Q A thin half ring of radius $35$ cm is uniformly charged with a total charge of $Q$ coulomb. If the magnitude of the electric field at centre of the half ring is $100$ V/m, then the value of $Q$ is \underline{\hspace{1.5cm}} nC. ($\epsilon_0=8.85\times10^{-12}\ \mathrm{C^2/Nm^2}$ and $\pi=3.14$)\etag{JEE Main 2026, 6 Apr, Mor.}
\optsfour{$2.14$}{$2.44$}{$3.25$}{$0.7$}

\Q $\sigma$ is the uniform surface charge density of a thin spherical shell of radius $R$. The electric field at any point on the surface of the spherical shell is:\etag{JEE Main 2024, 6 Apr, Mor.}
\optsfour{$\sigma/\epsilon_0R$}{$\sigma/\epsilon_0$}{$\sigma/2\epsilon_0$}{$\sigma/4\epsilon_0$}

\Qf A spherically symmetric charge distribution is considered with charge density varying as $\rho(r)=\rho_0\left(\dfrac34-\dfrac rR\right)$ for $r\le R$ and $\rho(r)=0$ for $r>R$. Here $r\ (r<R)$ is the distance from the centre O (as shown in figure). The electric field at point P will be:\etag{JEE Main 2022, 29 Jul, Mor.}
\fig{f94}
\optstwo{$\dfrac{\rho_0r}{4\epsilon_0}\left(\dfrac34-\dfrac rR\right)$}{$\dfrac{\rho_0r}{3\epsilon_0}\left(\dfrac34-\dfrac rR\right)$}{$\dfrac{\rho_0r}{4\epsilon_0}\left(1-\dfrac rR\right)$}{$\dfrac{\rho_0r}{5\epsilon_0}\left(1-\dfrac rR\right)$}

\Qf Consider a sphere of radius $R$ which carries a uniform charge density $\rho$. If a sphere of radius $\frac{R}{2}$ is carved out of it, as shown, the ratio $\left|\dfrac{\vec E_A}{\vec E_B}\right|$ of magnitude of electric field $\vec E_A$ and $\vec E_B$, respectively, at points A and B due to the remaining portion is:\etag{JEE Main 2020, 9 Jan, Mor.}
\fig{f147}
\optstwo{$\dfrac{17}{54}$}{$\dfrac{18}{54}$}{$\dfrac{18}{34}$}{$\dfrac{21}{34}$}

\Q For a uniformly charged ring of radius $R$, the electric field on its axis has the largest magnitude at a distance $h$ from its center. Then value of $h$ is:\etag{JEE Main 2019, 9 Jan, Mor.}
\optsfour{$\dfrac{R}{\sqrt5}$}{$\dfrac{R}{\sqrt2}$}{$R$}{$R\sqrt2$}

\Qf A thin semi-circular ring of radius $r$ has a positive charge $q$ distributed uniformly over it. The net field $\vec E$ at the centre $O$ is\etag{AIEEE 2010}
\fig{f220}
\optstwo{$\dfrac{q}{4\pi^2\epsilon_0r^2}\hat j$}{$-\dfrac{q}{4\pi^2\epsilon_0r^2}\hat j$}{$-\dfrac{q}{2\pi^2\epsilon_0r^2}\hat j$}{$\dfrac{q}{2\pi^2\epsilon_0r^2}\hat j$}

\Q A point charge $+q$ is placed at the origin. A second point charge $+9q$ is placed at $(d,0,0)$ in Cartesian coordinate system. The point in between them where the electric field vanishes is:\etag{JEE Main 2025, 2 Apr, Mor.}
\optstwo{$(3d/4,0,0)$}{$(d/4,0,0)$}{$(4d/3,0,0)$}{$(d/3,0,0)$}

\Q Two charges $q$ and $3q$ are separated by a distance $r$ in air. At a distance $x$ from charge $q$, the resultant electric field is zero. The value of $x$ is:\etag{JEE Main 2024, 31 Jan, Mor.}
\optstwo{$\dfrac{r}{3(1+\sqrt3)}$}{$\dfrac{(1+\sqrt3)}{r}$}{$\dfrac{r}{(1+\sqrt3)}$}{$r(1+\sqrt3)$}

\Q Two point charges A and B of magnitude $+8\times10^{-6}$ C and $-8\times10^{-6}$ C respectively are placed at a distance $d$ apart. The electric field at the middle point O between the charges is $6.4\times10^{4}\ \mathrm{N\,C^{-1}}$. The distance $d$ between the point charges A and B is:\etag{JEE Main 2022, 28 Jun, Eve.}
\optsfour{$2.0$ m}{$3.0$ m}{$1.0$ m}{$4.0$ m}

\Qf Figure shows a rod AB, which is bent in a $120^\circ$ circular arc of radius $R$. A charge $(-Q)$ is uniformly distributed over rod AB. What is the electric field $\vec E$ at the centre of curvature O?\etag{JEE Main 2021, 27 Aug, Eve.}
\fig{f117}
\optstwo{$\dfrac{3\sqrt3Q}{8\pi\epsilon_0R^2}(\hat i)$}{$\dfrac{3\sqrt3Q}{8\pi^2\epsilon_0R^2}(\hat i)$}{$\dfrac{3\sqrt3Q}{16\pi^2\epsilon_0R^2}(\hat i)$}{$\dfrac{3\sqrt3Q}{8\pi^2\epsilon_0R^2}(-\hat i)$}

\Qf Find the electric field at point P (as shown in figure) on the perpendicular bisector of a uniformly charged thin wire of length $L$ carrying a charge $Q$. The distance of the point P from the centre of the wire is $a=\dfrac{\sqrt3}{2}L$.\etag{JEE Main 2021, 26 Feb, Mor.}
\fig{f130}
\optstwo{$\dfrac{Q}{4\pi\epsilon_0L^2}$}{$\dfrac{Q}{3\pi\epsilon_0L^2}$}{$\dfrac{Q}{2\sqrt3\pi\epsilon_0L^2}$}{$\dfrac{\sqrt3Q}{4\pi\epsilon_0L^2}$}
\end{multicols}

\clearpage
\begin{center}{\Large\bfseries SELF-ASSESSMENT -- PRACTICE SHEET 1}\end{center}
\medskip
\renewcommand{\arraystretch}{1.9}
\noindent\begin{tabularx}{\linewidth}{|X|X|}\hline
\textbf{Time Allowed:} 60 Mins & \textbf{Time Taken:} \underline{\hspace{2.4cm}} Mins\\ \hline
\textbf{Total Questions:} 20 & \textbf{Questions Attempted:} \underline{\hspace{2cm}}\\ \hline
\textbf{Correct:} \underline{\hspace{2.6cm}} & \textbf{Incorrect:} \underline{\hspace{2.6cm}}\\ \hline
\textbf{Unattempted:} \underline{\hspace{2.2cm}} & \textbf{Accuracy:} \underline{\hspace{1.6cm}}\,\% \ {\footnotesize(Correct$\,/\,$Attempted$\,\times100$)}\\ \hline
\end{tabularx}
\renewcommand{\arraystretch}{1.0}
\bigskip
\textbf{Question-wise record} \ {\small(tick one: C = Correct, W = Wrong, S = Skipped)}\par\smallskip
{\small\renewcommand{\arraystretch}{1.55}
\noindent\begin{tabularx}{\linewidth}{|*{4}{>{\centering\arraybackslash}X|}}\hline
\textbf{Q.1}\ $\square$C\ $\square$W\ $\square$S & \textbf{Q.2}\ $\square$C\ $\square$W\ $\square$S & \textbf{Q.3}\ $\square$C\ $\square$W\ $\square$S & \textbf{Q.4}\ $\square$C\ $\square$W\ $\square$S \\ \hline
\textbf{Q.5}\ $\square$C\ $\square$W\ $\square$S & \textbf{Q.6}\ $\square$C\ $\square$W\ $\square$S & \textbf{Q.7}\ $\square$C\ $\square$W\ $\square$S & \textbf{Q.8}\ $\square$C\ $\square$W\ $\square$S \\ \hline
\textbf{Q.9}\ $\square$C\ $\square$W\ $\square$S & \textbf{Q.10}\ $\square$C\ $\square$W\ $\square$S & \textbf{Q.11}\ $\square$C\ $\square$W\ $\square$S & \textbf{Q.12}\ $\square$C\ $\square$W\ $\square$S \\ \hline
\textbf{Q.13}\ $\square$C\ $\square$W\ $\square$S & \textbf{Q.14}\ $\square$C\ $\square$W\ $\square$S & \textbf{Q.15}\ $\square$C\ $\square$W\ $\square$S & \textbf{Q.16}\ $\square$C\ $\square$W\ $\square$S \\ \hline
\textbf{Q.17}\ $\square$C\ $\square$W\ $\square$S & \textbf{Q.18}\ $\square$C\ $\square$W\ $\square$S & \textbf{Q.19}\ $\square$C\ $\square$W\ $\square$S & \textbf{Q.20}\ $\square$C\ $\square$W\ $\square$S \\ \hline
\end{tabularx}}
\bigskip
\textbf{Reason for mistakes} \ {\small(tick all that apply)}\par\smallskip
\begin{tabularx}{\linewidth}{@{}XXX@{}}$\square$ Concept not clear & $\square$ Formula forgotten & $\square$ Calculation mistake\\[3pt] $\square$ Misread the question & $\square$ Ran out of time & $\square$ Guess work\end{tabularx}
\bigskip
\textbf{Personal Feedback}\par\smallskip
\textit{What went well:}\par\vspace{2pt}
\noindent\rule{0pt}{1.35em}\hrulefill\\ \rule{0pt}{1.35em}\hrulefill\\[4pt]
\textit{Topics / questions I need to revise:}\par\vspace{2pt}
\noindent\rule{0pt}{1.35em}\hrulefill\\ \rule{0pt}{1.35em}\hrulefill\\[4pt]
\textit{My action plan before the next sheet:}\par\vspace{2pt}
\noindent\rule{0pt}{1.35em}\hrulefill\\ \rule{0pt}{1.35em}\hrulefill\\[8pt]
\textit{Teacher's remarks:}\par\vspace{2pt}
\noindent\rule{0pt}{1.35em}\hrulefill\\ \rule{0pt}{1.35em}\hrulefill\\[14pt]
Student's signature: \underline{\hspace{4cm}}\hfill Teacher's signature: \underline{\hspace{4cm}}
\end{document}
